\subsection{RL \& TOC-Based Simulation Model {[}REVISED DRAFT -- pending
author
review{]}}\label{rl-toc-based-simulation-model-revised-draft-pending-author-review}

\noindent\colorbox{yellow!40}{REVISED DRAFT --- pending author review. Kept alongside the original version of this section below; the two have not yet been reconciled (see the reviewer note at the end of this block).}\par\medskip

The TOC \& RL-based simulation model was also built upon the virtual
environment established in Section 4.1. The state space 𝒮 comprises 22
variables: two constraint-buffer variables --- the portion of
constraint-buffer time not yet arrived at the CCR, H\^{}CB\_wip,t, and
the portion already queued at the CCR, H\^{}CB\_queue,t --- together
with, for each product i, its work-in-process level, W\_i,t, and its
finished-goods (shipping-buffer) level, H\^{}SB\_i,t:

\emph{𝒮 = \{ H\^{}CB\_wip,t , H\^{}CB\_queue,t \} ∪ \{ W\_i,t ,
H\^{}SB\_i,t \textbar{} i = 1, \ldots, I \}, I = 10}

H\^{}CB\_wip,t and H\^{}CB\_queue,t are measured in hours of processing
time still owed to the CCR and take non-negative real values; W\_i,t and
H\^{}SB\_i,t are non-negative integers. Before reaching the policy
network, all state variables are normalized online using running mean
and variance statistics accumulated during training.

This state informs the agent, which independently selects, for each
product i, the number of units to release during the replenishment
interval, forming the action space:

\emph{𝒜(𝒮) = \{ M\_i,t \textbar{} m = 0, 1, \ldots, 60 \}, i = 1,
\ldots, I = 10}

This decision is made autonomously and is not explicitly modeled by
Equation 17. Because the I = 10 products are chosen independently and
simultaneously, 𝒜(𝒮) is a multi-discrete space with 10 sub-actions of 61
categories each --- 610 discrete outputs in the policy's action layer.
The upper bound M\_i,t = 60 units was set based on the highest order
size observed in the TOC-based model over {[}xx{]} replications
{[}confirm the scaling factor applied{]}. The production sequence
follows Equation 15, and, as in the TOC-based model, it is assumed that
(ρ\_i − c\_i) = 1 when defining τ\_i in Equation 10, since the original
work by Atwater and Chakravorty (2002) does not take into account the
revenue and variable costs of the products sold. The release of orders
follows Equation 19. These relationships are illustrated in Figure 5.

The agent receives a numerical reward R\_t+1 as a result of its action
(Sutton and Barto, 2018), composed of a fixed baseline, a throughput
term, and two penalty subfunctions:

\emph{(20) R\_t+1 = R\_0 + TH\_t+1 − P\^{}LS\_t+1 − P\^{}W\_t+1 −
P\^{}SB\_t+1}

R\_0 is a fixed baseline reward added at every step. TH\_t+1 is the
throughput term; because the original work by Atwater and Chakravorty
(2002) does not provide the information needed to compute throughput as
sold value, TH\_t+1 is set to zero during training, and its role is
absorbed entirely into the lost-sales penalty P\^{}LS\_t+1. Whenever
demand for product i exceeds its current finished-goods level
H\^{}SB\_i,t, the system immediately fails to make a sale and cannot
convert inventory into throughput (cf., Reference):

\emph{(21) P\^{}LS\_t+1 = λ\_LS · Σ(i=1..I) L\_i,t}

where L\_i,t is the number of units of product i lost to unmet demand
during interval t, and λ\_LS is a penalty weight (λ\_LS = 10 during
training).

The second measure is inventory, represented by the total quantity of
items rather than monetary values, for the same reason as throughput.
Inventory is penalized through two subfunctions, one for work-in-process
and one for finished goods, each summed across the I = 10 products:

\emph{(22) P\^{}W\_t+1 = λ\_W · Σ(i=1..I) W\_i,t}

\emph{(23) P\^{}SB\_t+1 = λ\_SB · Σ(i=1..I) H\^{}SB\_i,t}

with λ\_W = 0.002 and λ\_SB = 0.001. The constraint-buffer variables
(H\^{}CB\_wip,t, H\^{}CB\_queue,t) are part of the state 𝒮 but are not
penalized directly in the reward. The overall goal is to prevent lost
sales while maintaining the lowest possible inventory level (cf.,
Schragenheim, 2010b); accordingly, λ\_SB is set lower than λ\_W.

In the RL \& TOC-based simulation model, each decision step advances the
simulation clock by the replenishment interval τ = 168 hours, matching
the order-planning interval defined for the TOC-based model in Section
4.2. An episode begins at the start of a simulation run and proceeds
through successive decision steps until either (i) the reward R\_t+1
computed for a step turns negative, ending the episode immediately (cf.
Reference), or (ii) the simulation clock reaches the 200,000-hour
horizon, truncating the episode. By incorporating the reward function
(Equation 20) into the return function (Equation 2) and assuming a
discount factor γ = 0.99, along with rewards received at future time
steps (t′ \textgreater{} t), we construct the state-value function
(Equation 3), from which the agent is trained to seek the optimal policy
π*.

The algorithm used to train the optimal policy is PPO (Schulman et al.,
2017), implemented with the Stable-Baselines3 library (Raffin et al.,
2021), using a multi-input policy network with two hidden layers of 128
units each, shared between the actor and critic branches, that
approximates both the policy and the state-value function (Russell and
Norvig, 2010; Schulman et al., 2017). Observations and rewards were
normalized online (zero mean, unit variance, clipped at ±10,000)
throughout training. Training used 8 parallel environment instances, a
rollout length of 4,096 steps per environment, mini-batches of 2,048
transitions, and 5 optimization epochs per rollout, with a learning rate
that decayed linearly from 10⁻³ to 10⁻⁴ (Equation 7) over a maximum of
16,800,000 timesteps and 277,777 episodes, the virtual environment being
reset each time the step reward turned negative (cf. Reference). The
maximum number of episodes was set based on {[}define the reason{]}
(cf., Reference). Two criteria were used to terminate the training
process (Reference): (i) the maximum number of timesteps was reached, or
(ii) return saturation occurred, where the return remained stable over
{[}xx{]} episodes.

From a computational perspective, the training process was completed in
2 hours, utilizing 15 machines running in parallel, each with the
following specifications: an i5-13450HX processor, 16 GB of RAM, and an
NVIDIA GeForce RTX 4050 graphics card. A state-value of 0.996 was
achieved after 13,888 episodes, as shown in Figure 6a. The behavior of
the penalty subfunctions is also illustrated in Figure 6b--d.

\emph{{[}Reviewer note: Figure 6's original caption lists four
penalty-subfunction panels --- lost sales, work-in-process, CB level,
and SB level ((b)--(e)) --- but the reward above has only three: lost
sales, work-in-process, and SB/finished-goods (Equations 21--23), with
no CB term. Panel (d), the CB-level curve, should be removed, and the
current panel (e) (SB) should become the new panel (d), so the caption
reads: (a) return; (b) lost-sales penalty; (c) work-in-process penalty;
(d) finished-goods (SB) penalty. With this change, the body text's
existing ``Figure 6b--d'' citation is already correct and needs no
further edit.{]}}

\bigskip\noindent\hrulefill\par\smallskip\noindent\textbf{Original version} (kept for comparison --- the revised draft above and this original section have not yet been reconciled by the author).\par\bigskip
